% Work/MAT, 1. rocnik: referencni list k hodine 27.
% Skolni ceske nazvy; nikoli nova rectina. Srovnano s prehledem SLU AV CR.
% https://www.slu.cas.cz/sites/default/public_files/media/documents/2024/04/24/08/48/00/tabulka-alfabeta.pdf
\documentclass[11pt,a4paper]{article}
% Spolecny styl pro ucitelske pripravy MAT (Work/MAT)
% Verze 1.1, 2026-10-23
\NeedsTeXFormat{LaTeX2e}
\ProvidesPackage{priprava}[2026/10/23 v1.1 Sablona ucitelskych priprav MAT]

% Zakladni typografie a cestina
\RequirePackage[T1]{fontenc}
\RequirePackage[utf8]{inputenc}
\RequirePackage[czech]{babel}
\RequirePackage{lmodern}
\RequirePackage{microtype}
\RequirePackage{amsmath,amssymb}
\RequirePackage{xcolor}
\RequirePackage{enumitem}
\RequirePackage[most]{tcolorbox}
\RequirePackage{tabularx,booktabs,array}
\RequirePackage{fancyhdr}
\RequirePackage{needspace}
\RequirePackage{geometry}
\RequirePackage[hidelinks]{hyperref}

\geometry{a4paper,top=17mm,bottom=18mm,left=23mm,right=23mm}
\definecolor{notegray}{gray}{0.39}
\setlength{\parindent}{0pt}
\setlength{\parskip}{2.5pt plus 1pt minus 0.5pt}
\setlist[itemize]{topsep=2pt,itemsep=1pt,parsep=0pt,leftmargin=1.45em}
\setlist[enumerate]{topsep=2pt,itemsep=2pt,parsep=0pt,leftmargin=1.65em}
\setcounter{secnumdepth}{0}
\renewcommand{\arraystretch}{1.12}

% Zahlavi: vlevo rocnik | globalni tema, vpravo cislo: nazev hodiny
% Zapati zustava zcela prazdne.
\pagestyle{fancy}
\fancyhf{}
\renewcommand{\headrulewidth}{0.35pt}
\renewcommand{\footrulewidth}{0pt}
\setlength{\headheight}{14pt}
\newcommand{\nastavPripravu}[4]{%
  \fancyhead[L]{\small #1\enspace\textbar\enspace #2}%
  \fancyhead[R]{\small #3: #4}%
}

% Minimalisticke cernobile boxy. Kazda definice zvlast.
\tcbset{commonlesson/.style={enhanced,breakable=false,colback=white,colframe=black,
  colbacktitle=white,coltitle=black,boxrule=0.45pt,
  arc=0.8mm,outer arc=0.8mm,left=3mm,right=3mm,top=1.6mm,bottom=1.6mm,
  toptitle=1mm,bottomtitle=0.7mm,fonttitle=\bfseries,
  attach boxed title to top left={xshift=2.5mm,yshift=-1.15mm},
  boxed title style={colback=white,colframe=white,boxrule=0pt,left=1mm,right=1mm}}}
\newtcolorbox{zadaniBox}[1][]{commonlesson,title={Zadání},#1}
\newtcolorbox{prikladBox}[1][]{commonlesson,title={Příklad},#1}
\newtcolorbox{definiceBox}[1][]{commonlesson,title={Definice},#1}

% Spolecny box pro zadani a jeho reseni, bez mezery mezi castmi.
% Pouziti:
%   \begin{prikladReseniBox}
%     Zadani ulohy.
%     \tcblower
%     \begin{align*} ... \end{align*}
%   \end{prikladReseniBox}
% Spodni cast ma jemny sedy podklad a delici linku; lze v ni sazet
% vice radku matematiky. Obe casti maji jen jeden spolecny ramecek.
\tcbset{solutionlesson/.style={
  commonlesson,
  enhanced jigsaw,
  breakable,
  colbacklower=black!3,
  segmentation style={solid,draw=black!25,line width=0.35pt},
  fontlower=\small\color{notegray},
  before lower={\noindent\textbf{Řešení:}\enspace},
  leftlower=3mm,rightlower=3mm
}}
\newtcolorbox{prikladReseniBox}[1][]{solutionlesson,title={Příklad},#1}
\newtcolorbox{zadaniReseniBox}[1][]{solutionlesson,title={Zadání},#1}

% Pomucky pro autora pripravy: sedive poznamky; \reseni pro starsi pripravy
\newcommand{\poznamka}[1]{\par\begingroup\small\color{notegray}\itshape
  \textnormal{\textbf{Poznámka pro učitele.}}\enspace #1\par\endgroup}
\newcommand{\reseni}[1]{\par\begingroup\small\color{notegray}
  \textbf{Očekávané řešení:}\enspace #1\par\endgroup}
\newcommand{\blok}[1]{\needspace{5\baselineskip}\subsection*{#1}}
\newcommand{\vyrok}[1]{\uv{#1}}

\nastavPripravu{1. ročník}{Mocniny a odmocniny}{27}{Tabulka řeckých písmen}
\begin{document}
\blok{Řecká abeceda — 24 písmen}
V matematice používáme řecká písmena jako proměnné, konstanty i názvy veličin.
\textbf{S různými písmeny počítáme stejně jako s $a,b,x,y$.}
Následující české názvy odpovídají obvyklému školnímu čtení;
nejde o fonetický přepis současné řečtiny.
\vspace{2mm}
\begin{center}
\renewcommand{\arraystretch}{1.21}
\begin{tabularx}{\linewidth}{@{}c c l X@{}}
\toprule
\textbf{Velké} & \textbf{Malé} & \textbf{Čteme} & \textbf{Poznámka} \\\midrule
$\mathrm A$ & $\alpha$ & alfa & časté označení úhlu \\
$\mathrm B$ & $\beta$ & beta & časté označení úhlu \\
$\Gamma$ & $\gamma$ & gama & úhel, koeficient \\
$\Delta$ & $\delta$ & delta & $\Delta$ často značí změnu \\
$\mathrm E$ & $\varepsilon$ & epsílon & také tvar $\epsilon$ \\
$\mathrm Z$ & $\zeta$ & zéta &  \\
$\mathrm H$ & $\eta$ & éta &  \\
$\Theta$ & $\theta$ & théta & jiný tvar $\vartheta$ \\
$\mathrm I$ & $\iota$ & ióta &  \\
$\mathrm K$ & $\kappa$ & kappa &  \\
$\Lambda$ & $\lambda$ & lambda & například vlnová délka \\
$\mathrm M$ & $\mu$ & mí (mý) & například aritmetický průměr \\
$\mathrm N$ & $\nu$ & ný &  \\
$\Xi$ & $\xi$ & ksí &  \\
$\mathrm O$ & $o$ & omikron & tvarem jako latinské $o$ \\
$\Pi$ & $\pi$ & pí & $\pi\approx3{,}14159$ \\
$\mathrm P$ & $\rho$ & ró & například hustota \\
$\Sigma$ & $\sigma$ & sigma & $\sum$ je znak součtu; $\varsigma$ jiný tvar \\
$\mathrm T$ & $\tau$ & tau &  \\
$\Upsilon$ & $\upsilon$ & ypsilon &  \\
$\Phi$ & $\varphi$ & fí & další tvar $\phi$ \\
$\mathrm X$ & $\chi$ & chí &  \\
$\Psi$ & $\psi$ & psí &  \\
$\Omega$ & $\omega$ & omega & $\Omega$ může značit ohm \\\bottomrule
\end{tabularx}
\end{center}
\blok{Dvě miniukázky použití}
\[
\alpha^3\cdot\alpha^5=\alpha^8,
\qquad
\frac{\lambda^7\mu^4}{\lambda^2\mu}=\lambda^5\mu^3
\quad(\lambda,\mu\ne0).
\]
\textbf{Důležité:} $\alpha$ a $\beta$ jsou různá písmena. Nelze tedy
sečíst exponenty ve výrazu $\alpha^2\beta^3$ jako u stejného základu.
\small\textit{Poznámka: Velká řecká písmena často vypadají stejně jako
latinské kapitálky. U některých malých písmen existuje více tvarů.}\par\normalsize
\end{document}
